
\chapter{Token-Uniqueness Modes}
\label{ch:v154-token-uniqueness-modes}

The v1.5.3 patch made \(\TokUnique(\Pi)\) explicit. v1.5.4 keeps it as the active policy field, but separates two implementation modes so that the next machine sprint has a precise target.

\begin{definition}[Policy mode]
\label{def:v154-policy-mode}
In policy mode, \(\TokUnique(\Pi)\) is a field of the concrete package-token policy:
\[
  \TokIntro(\Pi,s,p),\quad \TokIntro(\Pi,t,p)
  \quad\Longrightarrow\quad
  \hsame(s,t).
\]
All v1.5.4 concrete Globalize exactness theorems may use this field as an assumption.
\end{definition}

\begin{definition}[Canonical-constructor mode]
\label{def:v154-canonical-mode}
In canonical-constructor mode, token introduction is implemented by a canonical token constructor, written schematically as
\[
  \TokCan(\Pi,s,p).
\]
The implementation must provide the no-confusion rule
\[
  \TokCan(\Pi,s,p),\quad \TokCan(\Pi,t,p)
  \quad\Longrightarrow\quad
  \hsame(s,t).
\]
If \(\TokIntro\) is defined by \(\TokCan\), then \(\TokUnique(\Pi)\) follows.
\end{definition}

\begin{lemma}[Canonical mode implies token uniqueness]
\label{lem:v154-canonical-implies-unique}
Assume \(\TokIntro(\Pi,s,p)\) is definitionally implemented by \(\TokCan(\Pi,s,p)\), and assume the no-confusion rule of Definition~\ref{def:v154-canonical-mode}. Then \(\TokUnique(\Pi)\) holds.
\end{lemma}

\begin{proof}
Given \(\TokIntro(\Pi,s,p)\) and \(\TokIntro(\Pi,t,p)\), unfold the implementation to obtain \(\TokCan(\Pi,s,p)\) and \(\TokCan(\Pi,t,p)\). The no-confusion rule gives \(\hsame(s,t)\), which is exactly \(\TokUnique(\Pi)\).
\end{proof}

\begin{principle}[Active mode for v1.5.4]
\label{prin:v154-active-token-mode}
The active paper theorem uses policy mode. Canonical-constructor mode is an implementation target, not a paper-level replacement, until a checked build records the implementation and no-confusion theorem.
\end{principle}

\begin{corollary}[Stable package reflection contract]
\label{cor:v154-stable-reflection-contract}
The package reflection theorem may be used in later v1.5.x patches only in one of two forms:
\begin{enumerate}
  \item base-relation reflection under \(\TokUnique(\Pi)\);
  \item base-relation reflection under a checked canonical-constructor theorem implying \(\TokUnique(\Pi)\).
\end{enumerate}
It may not be used for \(\psame_{\Pi}^{\mathrm{eq}}\) unless \(\ClosureReflect(\Pi)\) is separately proved.
\end{corollary}

\begin{proof}
The first form is Theorem~\ref{thm:v154-base-reflection-active}. The second form follows from Lemma~\ref{lem:v154-canonical-implies-unique}. The final exclusion is Corollary~\ref{cor:v154-eq-not-exported}.
\end{proof}
